\documentclass[10pt,a4paper]{amsart}
\usepackage[margin=19mm]{geometry}
\usepackage{amsmath,amssymb,mathtools}
\usepackage[scr=rsfs,cal=euler]{mathalfa}
\usepackage{xfrac}
\usepackage[unicode,hidelinks,bookmarks=false]{hyperref}
\newcommand{\ie}{i.e.\ }
\newcommand{\cf}{cf.\ }
\newcommand{\resp}{respectively\ }
\newcommand{\Subsection}{\subsection}
\providecommand{\nobreakdash}{\nobreak\hbox{-}}
\setlength{\emergencystretch}{2em}
\multlinegap0pt
\usepackage{bm}
\usepackage{bbm}
\usepackage{dsfont}
\usepackage{xspace}
\usepackage{cancel}
\usepackage{mathtools}
\usepackage{amsthm}
\makeatletter
\@ifundefined{theorem}{\newtheorem{theorem}{Theorem}}{}
\@ifundefined{proposition}{\newtheorem{proposition}[theorem]{Proposition}}{}
\makeatother
\DeclareMathOperator{\CE}{CE}
\DeclareMathOperator{\DR}{DR}
\DeclareMathOperator{\Der}{Der}
\DeclareMathOperator{\Sym}{Sym}
\newcommand{\bas}{\mathrm{bas}}
\newcommand{\bbT}{\mathbb T}
\newcommand{\cF}{\mathcal F}
\newcommand{\from}{\colon}
\newcommand{\gm}{\mathrm{gm}}
\newcommand{\g}{\mathfrak g}
\newcommand{\kk}{k}
\newcommand{\uuder}{\mathfrak u_{\mathrm{der}}}
\title[Transversal unfoldings of derived foliations]{Transversal unfoldings of derived foliations}
\author{Maur\'icio Corr\^ea}
\date{Source: arXiv:2606.04253v2 (CC BY 4.0)}
\begin{document}
\maketitle

\noindent\emph{Source and licence.} Transversal unfoldings of derived foliations. Version arXiv:2606.04253v2 (CC BY 4.0), distributed under the Creative Commons Attribution 4.0 International licence, as recorded in the arXiv licence metadata for this exact version and shown on the abstract page https://arxiv.org/abs/2606.04253v2. Full rights notice: licenses/arxiv-2606.04253-CC-BY-4.0.txt.\par\smallskip

\noindent\textit{Editorial scope.} Counted unit: the Proposition whose source label is \texttt{prop:intrinsic-comparison} in the article (source0.tex lines 1298--1308, source label \texttt{prop:intrinsic-comparison}), delivered together with the article's own complete original proof of it (source0.tex lines 1310--1374, one \texttt{proof} environment of 65 lines). No mathematics is written, repaired or re-derived: statement and proof are the article's own text line for line. Required context retained uncounted: none. Every definition, display and label of those lines is retained unchanged. Omitted from this package and not redistributed: 1--1297, 1309--1309, 1375--2529. Nothing inside a retained excerpt is omitted or altered: every retained body is delivered exactly as the article's lines are written, so no omission receipt of any kind accompanies this package. Citations: the retained excerpts carry no citation command at all, so no bibliography entry of the article is transcribed and no reference is redistributed. Reference adaptation: none was needed and none was made; every reference inside the delivered unit resolves inside it, so no unresolved reference or citation is produced. Printed slips kept literal: no printed word, symbol, hypothesis or number of the counted statement or of its proof was altered. Layout and class: the article's own class is \texttt{amsart} with packages including inputenc, fontenc, babel, geometry, microtype, mathpazo, euler, amsmath, amssymb, amsthm, mathtools, mathrsfs, bm, eucal; the wrapper is the lane's pinned \texttt{amsart} editorial wrapper at 10pt on A4, which restates the theorem-like environments the retained text needs and adds the article's own macro definitions that the retained text needs (\texttt{\textbackslash CE}, \texttt{\textbackslash DR}, \texttt{\textbackslash Der}, \texttt{\textbackslash Sym}, \texttt{\textbackslash bas}, \texttt{\textbackslash bbT}, \texttt{\textbackslash cF}, \texttt{\textbackslash from}, \texttt{\textbackslash g}, \texttt{\textbackslash gm}, \texttt{\textbackslash kk}, \texttt{\textbackslash uuder}), with extra wrapper packages (bm, bbm, dsfont, xspace, cancel, mathtools). The delivered numbering is the unit's own. Assets: no retained span contains a figure environment, a graphics-inclusion command, a PDF-page inclusion, a table or a TikZ picture; no image file is copied into this package, the package declares no asset dependency and no image file is redistributed with it. Licence: version arXiv:2606.04253v2 of this article is distributed under the Creative Commons Attribution 4.0 International licence, as recorded in the arXiv licence metadata for this exact version and shown on the abstract page https://arxiv.org/abs/2606.04253v2. A full rights notice is delivered at licenses/arxiv-2606.04253-CC-BY-4.0.txt. Characters: every retained excerpt is pure ASCII, so the delivered engine, XeLaTeX with its default Latin Modern text font, reports no missing character.
\begin{proposition}\label{prop:intrinsic-comparison}
Let \(B\to A\) be a smooth morphism between smooth commutative \(\kk\)-algebras, and let \(\cF_{A/B}\) be presented by a cofibrant strictly perfect relative dg-Lie algebroid
\[
 \rho\from\g\longrightarrow T_{A/B}
\]
through an equivalence of graded mixed cdga's
\[
 \DR(\cF_{A/B})\simeq\CE^*(\g).
\]
Then the dg-Lie algebra of basic weight-zero graded-mixed derivations of \(\CE^*(\g)\) identifies with \(D^1_{\bas}(\g/B)\). Under this identification the intrinsic inner action of \(\bbT_{\cF_{A/B}/B}\) is the map \(\iota\from\g\to D^1_{\bas}(\g/B)\). Hence the intrinsic controller is represented by the crossed module \([\g\to D^1_{\bas}(\g/B)]\), and by the ordinary quotient \(\uuder(\g/B)\) in the effective case.
\end{proposition}

\begin{proof}
We include the calculation because it identifies the graded-mixed formulation with the Lie--Rinehart formulation.

Set
 $
 R:=\CE^*(\g)=\Sym_A(\g^\vee[1]).
$
Throughout this calculation we adopt the convention that \(\CE^*(\g)=\Sym_A(\g^\vee[1])\), the mixed differential has bidegree \((-1,+1)\), and derivations are cohomological derivations of degree \(n\). The algebra \(R\) is free as a graded commutative algebra over \(A\), generated in weight one by \(\g^\vee[1]\). Hence a homogeneous weight-zero graded derivation \(\Theta\) of degree \(n\) of the underlying graded cdga is determined by the following two components:
\[
 \Theta_0:=\Theta|_A\from A\longrightarrow A[n]
\]
and
\[
 \Theta_1:=\Theta|_{\g^\vee[1]}\from
 \g^\vee[1]\longrightarrow \g^\vee[1][n],
\]
subject only to the graded Leibniz rule. The first component is a derivation \(\theta\in T_{A/\kk}^n\). Since \(\g\) is perfect, the second component may be dualised and desuspended; it is equivalently a \(\kk\)-linear operator
\[
 \delta\from\g\longrightarrow\g[n]
\]
satisfying the first-order Leibniz rule
\[
 \delta(ax)=\theta(a)x+(-1)^{n|a|}a\delta(x)
\]
for homogeneous \(a\in A\) and \(x\in\g\). Thus, before imposing compatibility with the mixed differential, a homogeneous weight-zero graded derivation of \(R\) is equivalent to a pair \((\theta,\delta)\), where \(\delta\) is a first-order operator on \(\g\) with symbol \(\theta\).

The graded mixed algebra \(R\) has two differentials. The internal cohomological differential is induced by the differential of the dg-Lie algebroid \(\g\). The mixed Chevalley--Eilenberg differential, denoted \(\epsilon_{CE}\), is the sum of the component dual to the anchor and the component dual to the bracket. A graded mixed derivation is a graded derivation commuting with \(\epsilon_{CE}\); the cohomological differential on the derivation complex is the commutator with the internal differential. Thus the structural condition is
\[
 [\Theta,\epsilon_{CE}]=0,
\]
not the assertion that \(\Theta\) is a closed element of the total derivation complex. We record the two substantive components of this equation.

First, evaluate the commutator on functions. For \(a\in A\), the anchor component of \(\epsilon_{CE}\) sends \(a\) to the functional \(x\mapsto \rho(x)(a)\). The identity \([\Theta,\epsilon_{CE}](a)=0\) says that, for every \(x\in\g\),
\[
 \theta(\rho(x)(a))-
\rho(\delta x)(a)-(-1)^{|\Theta||x|}\rho(x)(\theta(a))=0.
\]
Equivalently,
\[
 \rho(\delta x)=[\theta,\rho(x)] .
\]
Thus the symbol \(\theta\) and the operator \(\delta\) are compatible with the anchor.

Second, evaluate the commutator on a linear generator \(\alpha\in\g^\vee\). The bracket part of \(\epsilon_{CE}\alpha\) is dual to the Lie bracket on \(\g\). The identity \([\Theta,\epsilon_{CE}](\alpha)=0\), evaluated on homogeneous \(x,y\in\g\), becomes
\[
 \delta([x,y])=[\delta x,y]+(-1)^{|\delta||x|}[x,\delta y].
\]
Thus \(\delta\) is a derivation of the dg-Lie bracket. The internal differential of \(\g\) is not an additional algebraic constraint on the underlying graded derivation; it is the differential of the derivation complex, sending \((\theta,\delta)\) to the graded commutator with the internal differential. Under the correspondence above, this is the differential of \(D^1_{A/\kk}(\g)\).

Consequently graded-mixed derivations of \(\CE^*(\g)\) are pairs \((\theta,\delta)\) satisfying the structural conditions in the definition of \(D^1_{A/\kk}(\g)\), with the same cohomological differential. The basic condition is also read from the transverse symbol: \(\Theta\) is basic when the image of \(\theta\) under
\[
 T_{A/\kk}\longrightarrow A\otimes_B T_{B/\kk}
\]
lies in \(1\otimes T_{B/\kk}\). Together with the chosen base vector field \(\xi\), this identifies a basic graded-mixed derivation with the triple \((\theta,\delta,\xi)\) defining \(D^1_{\bas}(\g/B)\). Hence
\[
 \Der^{\gm}_{\bas}(\CE^*(\g))\cong D^1_{\bas}(\g/B)
\]
as dg-Lie algebras. The bracket is the commutator on both sides: on the CE side it is the commutator of derivations of the graded mixed cdga, and on the Lie--Rinehart side it is the commutator of first-order Lie derivations.

It remains to compare the inner actions. A tangent element \(x\in\g\) acts on \(R\) by the usual Cartan Lie derivative \(L_x=[\epsilon_{CE},\iota_x]\). Under the preceding description this derivation has symbol \(\rho(x)\) and acts on \(\g\) by \(\operatorname{ad}_x\). Hence it corresponds exactly to
\[
 \iota(x)=(\rho(x),\operatorname{ad}_x,0)\in D^1_{\bas}(\g/B).
\]
Taking the homotopy quotient by these inner derivations gives the crossed module \([\g\to D^1_{\bas}(\g/B)]\). If \(\iota\) is injective and is represented by a cofibration which is a degreewise monomorphism, the ordinary quotient \(\uuder(\g/B)\) is a strict model of the corresponding homotopy cofibre.
\end{proof}

\end{document}
